\section{Proofs of Theoretical Results}
\label{sec:appendix_proofs}

\subsection{Residual Cross-Term Bound}

\noindent\textbf{Lemma S1 (Residual cross-term bound).}
\label{lem:residual_gradient}
Let $\delta_i = \psi_i - \hat{\psi}_i^0$ with $\delta_0 \equiv 0$. On the simultaneous residual event in Assumption~2 of the main article, suppose a valid pointwise derivative bound $\|\nabla_x\delta_{i-1}(x)\|_2\leq G_\delta^{(i-1)}(x)$ is available on $\calX_{\mathrm{op}}$. Then
\begin{equation}
|L_{\Delta\hat f}\delta_{i-1}(x)|
\leq G_\delta^{(i-1)}(x)\,\beta\|\sigmaGP(x)\|_2.
\label{eq:G_delta_cross_app}
\end{equation}

\begin{proof}
By Cauchy--Schwarz,
\[
|L_{\Delta\hat f}\delta_{i-1}|
=|\nabla_x\delta_{i-1}^{\top}\Delta\hat f|
\leq\|\nabla_x\delta_{i-1}\|_2\|\Delta\hat f\|_2.
\]
The simultaneous residual event gives $\|\Delta\hat f(x)\|_2\leq\beta\|\sigmaGP(x)\|_2$ over the modeled outputs, which yields Eq.~\eqref{eq:G_delta_cross_app}. Smoothness on a compact operating set ensures that finite global derivative bounds exist when the required derivatives are bounded, but it does not identify them with $\|\nabla\hat\psi_{i-1}^0\|_2$.
\end{proof}

The online code uses $\|\nabla\hat\psi_{i-1}^0\|_2$ as a commissioning proxy for $G_\delta^{(i-1)}$. Theorem~1 applies to that implementation only when Assumption~3 of the main article independently verifies that the resulting total margin dominates the true top-level HOCBF residual on $\calX_{\mathrm{op}}$.

\subsection{Full Proof of Theorem~1}

\begin{proof}[Proof of Theorem~1]
\textbf{Step 1 (Simultaneous residual event).} Assumption~2 of the main article states that $|\Deltaf_j(x)-\mu_{\mathrm{GP},j}(x)|\leq\beta\sigmaGPj{j}(x)$ holds simultaneously for the $q=3$ modeled residual outputs and all $x\in\calX_{\mathrm{op}}$ with probability at least $1-\delta$. Condition on this event. The commissioning multiplier used in the numerical studies does not establish this event by itself.

\textbf{Step 2 (Compositional bound).} We show $|\delta_i(x)| \leq \sigma_i(x)$ by induction on $i$.
\begin{itemize}
    \item \emph{Base case ($i=1$):} $\delta_1 = \sum_j (\partial h/\partial x_j) \Delta \hat{f}_j$. By Step~1 and triangle inequality, $|\delta_1| \leq \beta \sum_j |\partial h/\partial x_j| \sigmaGPj{j} = \sigma_1$.
    \item \emph{Inductive step ($i \geq 2$):} Assume $|\delta_{i-1}| \leq \sigma_{i-1}$. The residual recursion in the main article gives $|\delta_i| \leq |L_{\Delta \hat{f}} \hat{\psi}_{i-1}^0| + (|L_{\fhat}|_{\mathrm{op}} + k_i)|\delta_{i-1}| + |L_{\Delta \hat{f}} \delta_{i-1}|$. The first term is bounded as in the base case. The second term is bounded by the induction hypothesis. Lemma~S1 bounds the third term when a valid $G_\delta^{(i-1)}$ is available. Assumption~3 requires the operator-norm and cross-term quantities used for the certificate to dominate these terms. Summing gives $|\delta_i| \leq \sigma_i$.
\end{itemize}
The implementation proxy is not substituted for $G_\delta^{(i-1)}$ in this proof unless Assumption~3 validates that substitution.

\textbf{Step 3 (Total perturbation).} Combining the inductive bounds from Step~2 with the coupling-coefficient bound in Eq.~\eqref{eq:sigma_ctrl_app} yields $|\Delta_{\psi}(x,v)| \leq \sigma_m(x) + \sum_{j=1}^{m-1} c_j \sigma_j(x) + \sigma_{\mathrm{ctrl}}(x) = \sigmatotal(x)$ for every admissible deviation input $v$. For $\epsilonKappa = 1$, the applied margin is $\sigmatotal(x)$.

\textbf{Step 4 (Forward invariance).} For $v^*$ satisfying $A_0(x)v^* \leq b_0(x)-\sigmatotal(x)$: $\psi_m^{\mathrm{true}}(x,v^*) = \hat{\psi}_m^0(x,v^*) + \Delta_{\psi}(x,v^*) \geq \sigmatotal(x) - |\Delta_{\psi}(x,v^*)| \geq 0$. By the HOCBF forward-invariance theorem cited in the main article, the safe set $\calC$ is forward invariant.
\end{proof}

\subsection{Coupling Coefficient Gap Bound}

The coupling coefficient gap is bounded by:
\begin{equation}
    \sigma_{\mathrm{ctrl}}(x) = \beta \sum_{j} \left|\frac{\partial (L_g L_{\fhat}^{m-1} h)}{\partial x_j}\right| \sigmaGPj{j}(x) \cdot u_{\max}
    \label{eq:sigma_ctrl_app}
\end{equation}
where $u_{\max} = \sup_{u \in \calU} \|u\|$. For the CCS benchmark, $\sigma_{\mathrm{ctrl}} / \sigmatotal < 0.15$ for both constraints.
